\documentclass[11pt,a4paper]{article}
\usepackage[T1]{fontenc}\usepackage[utf8]{inputenc}\usepackage{lmodern}\usepackage{microtype}
\usepackage[a4paper,margin=0.82in,headheight=15pt,headsep=18pt,footskip=25pt]{geometry}
\usepackage{amsmath,amssymb,amsthm,mathtools,bm}\usepackage{graphicx}\usepackage{booktabs,tabularx,longtable,array,multirow}
\usepackage{xcolor}\usepackage{enumitem}\usepackage{fancyhdr}\usepackage{titlesec}\usepackage{caption}\usepackage{subcaption}\usepackage{hyperref}\usepackage[nameinlink,noabbrev]{cleveref}\usepackage[numbers,sort&compress]{natbib}\usepackage{float}\usepackage{placeins}\usepackage{listings}
\definecolor{ink}{HTML}{243447}\definecolor{teal}{HTML}{286B70}\definecolor{muted}{HTML}{64748B}\definecolor{pale}{HTML}{EEF3F5}
\hypersetup{colorlinks=true,linkcolor=teal,citecolor=teal,urlcolor=teal,pdfauthor={Devanik Debnath},pdftitle={Epsilon-Nets of Antigenic Space: Conditional Bounds for Broad-Spectrum Antibody Cocktails}}
\titleformat{\section}{\Large\bfseries\color{ink}}{\thesection}{.65em}{}\titleformat{\subsection}{\large\bfseries\color{teal}}{\thesubsection}{.6em}{}\titleformat{\subsubsection}{
\normalsize\bfseries\color{ink}}{\thesubsubsection}{.5em}{}
\titlespacing*{\section}{0pt}{2ex plus .5ex minus .2ex}{1ex}\titlespacing*{\subsection}{0pt}{1.5ex plus .4ex minus .2ex}{.7ex}
\setlength{\parskip}{.38em}\setlength{\parindent}{1.1em}\setlist{nosep,leftmargin=1.6em}
\captionsetup{font=small,labelfont=bf,justification=justified,singlelinecheck=false}
\pagestyle{fancy}\fancyhf{}\fancyhead[L]{\small\color{muted}Epsilon-Nets of Antigenic Space}\fancyhead[R]{\small\color{muted}Devanik Debnath · Preprint}\fancyfoot[C]{\small\thepage}\renewcommand{\headrulewidth}{.35pt}

\newtheorem{theorem}{Theorem}[section]
\newtheorem{proposition}[theorem]{Proposition}
\newtheorem{lemma}[theorem]{Lemma}
\newtheorem{corollary}[theorem]{Corollary}\theoremstyle{definition}
\newtheorem{definition}[theorem]{Definition}\theoremstyle{remark}
\newtheorem{remark}[theorem]{Remark}

\newcommand{\Acal}{\mathcal A}
\newcommand{\Xcal}{\mathcal X}
\newcommand{\Rcal}{\mathcal R}
\newcommand{\eps}{\varepsilon}
\newcommand{\VC}{\mathrm{VCdim}}
\newcommand{\IC}{\mathrm{IC}_{50}}
\newcommand{\Prb}{\mathbb P}
\newcommand{\E}{\mathbb E}
\newcommand{\R}{\mathbb R}
\newcommand{\one}{\mathbf 1}
\begin{document}
\hypersetup{pageanchor=false}
\begin{titlepage}\centering\vspace*{.45in}
{\large\color{teal}\textsc{Theoretical Research Preprint}}\\[1em]
{\huge\bfseries\color{ink}Epsilon-Nets of Antigenic Space}\\[.45em]
{\LARGE\bfseries Conditional Bounds for Broad-Spectrum Antibody Cocktails}\\[1em]
{\large Pfaffian and o-minimal range complexity, multi-covers, and the price of robust protection}\\[1.8em]
{\Large Devanik Debnath}\\[.3em]
{
\normalsize ECE '26 · National Institute of Technology Agartala, India}\\
{
\normalsize \href{https://github.com/Devanik21}{GitHub: Devanik21} · \href{https://orcid.org/0009-0001-8310-6921}{ORCID: 0009-0001-8310-6921}}\\[1.25em]
\begin{tabular}{rl}\textbf{Status:}& Independent theoretical preprint; not peer reviewed\\\textbf{Version:}&1.0.0\\\textbf{Date:}&1 November 2026\\\textbf{Evidence:}& Formal framework and synthetic verification only\end{tabular}
\vfill
\begin{minipage}{.93\textwidth}\small
\textbf{Abstract.} Broad-spectrum neutralization is often approached by optimizing antibody or immunogen combinations against finite panels of pathogen variants. The number of components required for reliable coverage remains uncertain in several important settings, including HIV broadly neutralizing antibody combinations and recombinant antivenoms. We formulate passive antibody selection as a range-space problem. Candidate antibodies form a measured ground set $(\Acal,\mu)$; each antigenic target $x$ defines a neutralization range $R_x=\{a\in\Acal:\nu(a,x;c)\geq\theta\}$. If every target range has candidate-repertoire mass at least $\eps$ and the family has finite VC dimension $d$, classical $\eps$-net theory guarantees a finite antibody subset intersecting every target range, with size $O((d/\eps)\log(1/\eps))$. Thus, under fixed complexity and minimum-mass assumptions, the sufficient bound has no explicit dependence on the number of target variants. We show how Hill-type response models can yield definable threshold ranges and explain why a two-dimensional antigenic map does not by itself set $d=2$. For multiple distinct neutralizers, we compare the published finite-uniform $\eps$-$k$-net theorem with a weighted small-atom sufficient construction proved here. A robust-neighborhood formulation makes the coverage loss under bounded escape perturbations explicit. A deterministic synthetic benchmark verifies the range geometry and finite multi-cover implementation; it is not a biological dataset or clinical validation. The contribution is a conditional mathematical bridge and an explicit assumption ledger, not a universal vaccine, a new general $\eps$-net theorem, or a solved immunological cocktail-size problem.
\end{minipage}
\vfill\begin{flushleft}\small\textbf{Keywords:} epsilon-nets; VC dimension; Pfaffian functions; o-minimality; neutralizing antibodies; antigenic cartography; multi-cover; robust optimization; active learning.\end{flushleft}
\end{titlepage}
\hypersetup{pageanchor=true}\pagenumbering{roman}
\begingroup\footnotesize\setlength{\parskip}{0pt}\tableofcontents\endgroup\clearpage
\listoffigures\listoftables\clearpage\pagenumbering{arabic}

\section{Problem statement and the meaning of diversity-free}
A broad-spectrum antibody cocktail is a finite set of molecular binders intended to neutralize a declared family of antigenic targets. The number of targets in a study panel is rarely the only difficulty. Targets differ in epitope exposure, glycosylation, conformational state, assay context, and sensitivity to individual antibodies. A cocktail can be broad in a sampled panel and still fail on an unmeasured or newly evolved variant. Conversely, a target panel can be very large while the family of candidate subsets that neutralize individual targets has low combinatorial complexity. These observations motivate a mathematical question: under what conditions can one select a finite set of antibody candidates that intersects every neutralization range above a prespecified coverage threshold?

The phrase \emph{diversity-free} is deliberately narrow. The proposed upper bound has no explicit factor equal to the cardinality $|\Xcal|$ of the finite target panel when the VC dimension of the range family is held fixed. It is not independent of biological diversity in every sense. Target diversity can lower the worst-case neutralizer mass $\eps$, increase the complexity parameter $d$, invalidate a model family, or introduce targets for which no candidate neutralizes. Each change can destroy the numerical usefulness of the bound. We therefore use the phrase only for independence from target-set cardinality under uniform complexity and margin assumptions.

A second distinction concerns the word \emph{minimal}. An $\eps$-net theorem is an existence theorem for a sufficient hitting set. It does not identify the exact minimum cocktail, does not guarantee that a greedy algorithm finds an optimum, and does not prove that a vaccine immunogen mixture induces the selected antibodies in a human host. Minimum-cardinality selection is a data-dependent hitting-set or set-multicover optimization problem. The theorem supplies a worst-case certificate that a finite solution exists under specified conditions.

The present target is a conditional theory of broad-spectrum antibody selection. It is potentially relevant to passive immunization and recombinant antivenom design, where antibody components themselves are selected. Translation to active vaccination requires a further map from immunogens and host context to the induced antibody distribution. The mathematical guarantee begins only after that induced distribution has been defined and its coverage margin established. Figure~\ref{fig:pipeline} summarizes the certificate path.
\begin{figure}[t]\centering\includegraphics[width=.84\linewidth]{figures/fig09_epsilon_pipeline.pdf}\caption{Conditional coverage pipeline. The repertoire measure, target-indexed neutralization ranges, and family complexity must be specified before an $\eps$-net theorem can be applied. The certificate covers only ranges with $\mu(R_x)\geq\eps$.}\label{fig:pipeline}\end{figure}

\section{Biomedical motivation and adjacent evidence}
In HIV research, broadly neutralizing antibodies (bnAbs) are evaluated against panels of genetically diverse Env pseudoviruses, and combinations are investigated to improve breadth and potency. Reviews describe the challenge of selecting specificities that jointly cover circulating viruses and limit escape; clinical and preclinical work has often compared cocktails of a few bnAbs against large neutralization panels \citep{JulgBarouch2021,Thavarajah2024}. A 1 November 2026 review reports progress in inducing several bnAb precursor lineages in humans, while identifying coordinated maturation of two to three specificities and durable serum activity as remaining vaccine-development tasks \citep{HaynesYehBaden2026}. These are active-immunogen problems, not equivalent to selecting a passive antibody mixture from an existing candidate library.

Recombinant antivenom research presents a parallel but distinct selection problem. A 2025 review describes mixtures of characterized monoclonal antibodies as a promising route for toxin neutralization and emphasizes the need to extend broad-neutralizing antibodies across toxin classes \citep{Menzies2025}. Toxin targets may form a structured collection of families, isoforms, species, and clinically relevant activities. A single low-dimensional Euclidean chart may be inappropriate; a range-space representation can nevertheless be defined over a discrete or structured target domain.

Antigenic cartography provides low-dimensional representations of serological relationships and has been applied to influenza strain characterization and vaccine-strain decisions \citep{Sitaras2020}. Such maps summarize measured antigenic distances. They do not imply that antibody candidates occupy the same two-dimensional metric space, nor that neutralization regions are Euclidean disks. Cartography is a possible feature construction, not proof of low VC dimension.

The CATNAP resource compiles published neutralization measurements alongside HIV Env sequences and exposes IC$_{50}$, IC$_{80}$, and ID$_{50}$ data for analysis \citep{Yoon2015}. Its download page listed assay, antibody, and virus tables for the 1 October 2026 snapshot when checked in this research cycle. The download endpoint could not be retrieved from the available execution environment; no CATNAP rows were downloaded, harmonized, or used below. CATNAP is a concrete next empirical resource, not an analyzed dataset in this release.

\begin{table}[t]\centering\small\caption{Selected adjacent literature and its relationship to this paper.}\begin{tabularx}{\linewidth}{p{.20\linewidth}p{.35\linewidth}X}\toprule Area&Established contribution&Boundary relative to present framework\\\midrule HIV bnAb combinations&Empirical panels compare breadth, potency, complementarity, and escape risk \citep{JulgBarouch2021,Thavarajah2024}.&No general distribution-free antibody-range selection bound follows automatically from those studies.\\ Human bnAb induction&Sequential immunogens can engage and mature distinct B-cell lineages \citep{HaynesYehBaden2026}.&Immunogens are not antibodies; host-dependent induction is a separate operator.\\ Recombinant antivenoms&Defined monoclonal mixtures are a translationally plausible strategy \citep{Menzies2025}.&Toxin-family heterogeneity demands a separate target domain and endpoint.\\ Antigenic cartography&Serological relationships can be embedded in low-dimensional maps \citep{Sitaras2020}.&Map dimension is not antibody-range VC dimension.\\ $\eps$-nets&Finite VC dimension yields finite hitting sets \citep{HausslerWelzl1987}.&The biological range mapping and lower-mass premise need independent validation.\\ $\eps$-$t$-nets&Published multi-coverage constructions extend classical nets \citep{AlonEtAl2022}.&The theorem has finite-uniform hypergraph hypotheses, not arbitrary weighted biological repertoires.\\\bottomrule\end{tabularx}\label{tab:literature}\end{table}

\section{Formal range-space model}
Let $\Acal$ be a set of candidate antibodies, antibody-derived binders, or other explicitly defined neutralizing agents. A probability measure $\mu$ on $\Acal$ represents the repertoire from which components are eligible to be selected. In an empirical study, $\mu$ might be uniform over a discovery library, weighted by availability, or defined by a pre-registered sampling distribution. Its meaning must be fixed before examining coverage outcomes; it cannot be reweighted after observing difficult targets.

Let $\Xcal$ denote a target domain. A target $x\in\Xcal$ may be a viral isolate, antigenic variant, toxin, toxin isoform, or structured target profile. Fix assay conditions $c$ (concentration, matrix, endpoint, incubation protocol, and other controlled factors) and a neutralization threshold $\theta\in(0,1)$. For a response function $\nu:\Acal\times\Xcal\times\mathcal C\to[0,1]$, define
\begin{equation}R_x(c,\theta)=\{a\in\Acal:\nu(a,x;c)\geq\theta\}.\label{eq:range}\end{equation}
The range space is $(\Acal,\Rcal)$, where $\Rcal=\{R_x(c,\theta):x\in\Xcal\}$. If two targets induce the same set, they are indistinguishable for this thresholded coverage question, even if their molecular sequences differ.

A finite subset $S\subseteq\Acal$ is an $\eps$-net if it intersects every range whose measure is at least $\eps$. The key premise is
\begin{equation}\inf_{x\in\Xcal}\mu(R_x(c,\theta))\geq\eps>0.\label{eq:margin}\end{equation}
Equation~\eqref{eq:margin} is a minimum-mass or breadth-margin assumption. It is not implied by finite VC dimension, by a compact-looking antigenic map, or by a large antibody library. A target for which all candidates fail has mass zero and cannot be guaranteed by an $\eps$-net. The domain must therefore be explicit, and excluded target strata must be reported.

\begin{definition}[VC dimension] A finite subset $B\subseteq\Acal$ is shattered by $\Rcal$ if $\{R\cap B:R\in\Rcal\}=2^B$. The VC dimension $d=\VC(\Acal,\Rcal)$ is the largest size of a shattered finite subset, or infinity if arbitrarily large finite subsets are shattered.\end{definition}

\begin{table}[t]\centering\small\caption{Notation and operational interpretations.}\begin{tabularx}{\linewidth}{p{.15\linewidth}X p{.30\linewidth}}\toprule Symbol&Definition&Operational interpretation\\\midrule $\Acal$&Candidate universe&Discovery library or repertoire support\\ $\mu$&Probability measure on $\Acal$&Predeclared eligibility or sampling weights\\ $\Xcal$&Target domain&Finite panel or specified continuum\\ $\nu(a,x;c)$&Neutralization response&Assay-specific response at conditions $c$\\ $\theta$&Response threshold&Prespecified functional endpoint\\ $R_x$&Candidates neutralizing target $x$&Target-indexed subset of antibody candidates\\ $\eps$&Lower bound on $\mu(R_x)$&Minimum candidate mass per target\\ $d$&$\VC(\Acal,\Rcal)$&Combinatorial complexity of target ranges\\ $k$&Required distinct hits&Redundancy demand per range\\ $\eta$&Largest atom of $\mu$&Maximum mass of one candidate\\ $\rho$&Escape-neighborhood radius&Declared local perturbation budget\\\bottomrule\end{tabularx}\label{tab:notation}\end{table}

\section{Hill response and threshold geometry}
A common dose-response idealization is a Hill curve:
\begin{equation}\nu(a,x;c)=\frac{c^{h(a,x)}}{c^{h(a,x)}+\IC(a,x)^{h(a,x)}}=\sigma\!\left(h(a,x)[\log c-\log\IC(a,x)]\right),\label{eq:hill}\end{equation}
where $c>0$ is concentration, $h(a,x)>0$ is a Hill slope, $\IC(a,x)>0$ is an IC$_{50}$ scale, and $\sigma(t)=(1+e^{-t})^{-1}$. Equation~\eqref{eq:hill} is a response family, not a claim that every assay follows an exact Hill law. Fitted curves can be affected by censoring, plate effects, incomplete neutralization, and between-study differences.

For fixed $c$ and threshold $\theta$, the condition $\nu\geq\theta$ is equivalent to
\begin{equation}\log c-\log\IC(a,x)\geq h(a,x)^{-1}\log\frac{\theta}{1-\theta}.\label{eq:hillthreshold}\end{equation}
At $\theta=1/2$, the right side is zero, and the threshold reduces to $\IC(a,x)\leq c$, independently of positive Hill slope. This simplification is used in the toy geometry. At other thresholds, the Hill slope affects range boundaries and must be part of the complexity audit.

A special case embeds candidates as points $z(a)\in\R^2$ and targets as centers $m(x)\in\R^2$, with
\begin{equation}\log\IC(a,x)=\kappa(\|z(a)-m(x)\|_2-r(x)),\qquad\kappa>0.\label{eq:icmodel}\end{equation}
At $c=1$ and $\theta=1/2$, the range is a variable-radius disk in candidate-coordinate space. The family of all disks in the plane has VC dimension three. This gives a transparent example, not a general inference from the dimension of an antigenic map.
\begin{figure}[t]\centering\includegraphics[width=.78\linewidth]{figures/fig01_hill_response.pdf}\caption{Synthetic Hill-response profiles with a common IC$_{50}$ radius and different slopes. At the 50\% threshold, the profiles induce the same boundary; at other thresholds their ranges differ.}\label{fig:hill}\end{figure}

\section{Definability and finite range complexity}
The mathematical bridge from a response formula to a combinatorial dimension uses uniform definability. Suppose candidates are represented by parameters $a\in A\subseteq\R^p$ and target profiles by $x\in X\subseteq\R^q$. Assume a fixed first-order formula $\varphi(a,x)$ in an o-minimal expansion of the ordered real field defines
\begin{equation}R_x=\{a\in A:\varphi(a,x)\}.\label{eq:definable}\end{equation}
Uniform definability means the same formula defines every target-specific range, with formula complexity fixed while $x$ varies. O-minimality constrains the geometry of definable sets and excludes arbitrary oscillatory behavior on the real line. Uniformly definable families in o-minimal structures have finite combinatorial complexity; parameter dimension and formula structure govern quantitative bounds \citep{JohnsonLaskowski2010,Wilkie1996,vdDMM1994}.

Hill-type models sit naturally near this setting. The exponential and logarithm are definable in the real field expanded by exponentiation; the logistic sigmoid is Pfaffian because $\sigma'=\sigma(1-\sigma)$. If potency and feature functions are built from a fixed finite combination of polynomial, exponential, logarithmic, and restricted analytic/Pfaffian operations on a domain where all expressions are valid, the threshold relation can be represented by a fixed definable formula. A formula with arbitrary lookup tables, uncontrolled sequence-specific branches, or unconstrained learned components need not preserve a useful complexity bound.

\begin{proposition}[Finite VC dimension of a uniformly definable family] Let $A\subseteq\R^p$ and $X\subseteq\R^q$ be definable in an o-minimal structure, and let one fixed formula $\varphi(a,x)$ define the family in \eqref{eq:definable}. Then $\Rcal=\{R_x:x\in X\}$ has finite VC dimension.\end{proposition}
\begin{proof} This is the standard finite-VC property of uniformly definable families in o-minimal structures. The proposition is qualitative: it does not supply a useful numerical value of $d$ without formula-specific complexity data.\end{proof}

Pfaffian theory can provide explicit bounds in terms of ambient dimension, chain length, polynomial degrees, and the format of a computational architecture. VC-dimension bounds for sigmoidal and Pfaffian neural networks, and for programs with transcendental operations, establish quantitative results for specified formats \citep{KarpinskiMacintyre1997,MontanaPardo2009}. Such constants cannot be replaced by the number of coordinates in an antigenic map. No numerical Pfaffian format certificate for actual HIV neutralization or toxin neutralization is derived here.

\begin{figure}[t]\centering\includegraphics[width=.78\linewidth]{figures/fig03_epsilon_scaling.pdf}\caption{Dimension-dependent scaling index $d\eps^{-1}\log(1/\eps)$ for $d=3$. The universal constant is omitted; this is not a calibrated cocktail-size estimate.}\label{fig:scaling}\end{figure}

\section{What Pfaffian structure does and does not buy}
A Pfaffian chain is a finite sequence of differentiable functions whose derivatives are polynomial expressions in ambient variables and earlier chain elements. This structure makes it possible to bound zeros and connected components of sets defined through sign conditions. For neutralization models, one may specify a chain containing feature transformations, potency surfaces, and response nonlinearities, then express a threshold boundary as a sign condition. The bound depends on the complete format tuple: number of variables, chain length, degrees, and number of sign tests.

There are two proof routes. A qualitative route shows that the response family is uniformly definable in an o-minimal structure and concludes finite VC dimension. A quantitative route encodes the classifier in a specified Pfaffian or computational formula format and applies a published bound appropriate to that format. The latter is harder because assay normalization, categorical epitope interactions, sequence-specific exceptions, and learned feature maps can change the formula class. A quantitative theorem for one architecture cannot simply be transferred to another model because both are described as Hill-like.

The tempting statement “two-dimensional antigenic coordinates imply $d=3$” is false in general. It is valid only when candidate antibodies are points in the plane and each target range is a Euclidean disk (or a subfamily of disks) in that candidate plane. If target coordinates are two-dimensional but binders occupy a high-dimensional sequence space, the family is not automatically a disk family. Conversely, high-dimensional candidate coordinates can still induce low-complexity ranges if the boundary family is restricted. Parameter dimension, formula complexity, and the induced subsets jointly determine VC dimension.

A data-driven neural predictor can be useful for ranking candidate-antigen pairs, but unconstrained neural architectures do not automatically have small or useful biological VC bounds. Even when a formal VC bound exists, it may be too loose to inform cocktail selection. Model selection must therefore compare prediction quality, calibration, formula complexity, and the resulting worst-case margin. Interpretability is not obtained merely by naming a model “Pfaffian.”

\begin{table}[t]\centering\small\caption{Assumptions required for mathematical transfer and representative failure modes.}\begin{tabularx}{\linewidth}{p{.19\linewidth}p{.35\linewidth}X}\toprule Assumption&Mathematical role&Failure mode\\\midrule Fixed candidate universe&Defines selected points&Candidates generated adaptively or removed post hoc\\ Prospective measure $\mu$&Defines range mass&Weights are reselected after observing coverage\\ Fixed assay context&Makes $R_x$ well-defined&Assay, matrix, temperature, concentration, or endpoint shifts\\ Uniform response formula&Supports definability and finite VC&Arbitrary branches or changing model architecture\\ Positive margin $\eps$&Every range is heavy&Resistant variants have zero or near-zero neutralizer mass\\ Single-agent incidence&Allows range intersection&Synergy, antagonism, competition, pharmacokinetics omitted\\\bottomrule\end{tabularx}\label{tab:assumptions}\end{table}

\section{The classical epsilon-net theorem in antibody space}
The central combinatorial theorem is classical. Haussler and Welzl introduced $\eps$-nets for abstract range spaces and established sampling bounds for finite VC dimension \citep{HausslerWelzl1987}. In the present orientation, antibody candidates are points and antigenic targets induce ranges. We seek a small set of antibodies that hits every heavy neutralization range, not a small set of antigens that is covered by every antibody.

\begin{theorem}[VC epsilon-net bound]\label{thm:epsilon} Let $(\Acal,\Rcal)$ be a range space with VC dimension $d\geq1$ and let $\mu$ be a probability measure on $\Acal$. For $0<\eps<1$, there exists an $\eps$-net $S\subseteq\Acal$ with
\begin{equation}|S|\leq C_{\rm HW}\frac{d}{\eps}\log\frac{1}{\eps},\label{eq:epsilonbound}\end{equation}
for a universal constant (with standard adjustments in constants and logarithms for finite or small-parameter regimes). If $\mu(R_x)\geq\eps$ for every declared target $x$, every such target range is intersected by $S$.\end{theorem}

The bound has no explicit $|\Xcal|$ term. If the response family retains dimension $d$ as the target panel grows and its lower-mass guarantee remains at least $\eps$, the sufficient sample size is controlled by $d$ and $\eps$, not the number of listed variants. This may be useful when a finite panel samples a larger antigenic continuum. It can also be uninformative if new target strata increase $d$ or reduce the margin.

Random samples provide a related high-probability guarantee. Standard VC sampling bounds give a sufficient sample size of order
\begin{equation}n=O\!\left(\frac{1}{\eps}\left[d\log\frac{1}{\eps}+\log\frac{1}{\delta}\right]\right),\label{eq:randomsample}\end{equation}
for failure probability at most $\delta$, up to universal constants and standard parameter qualifications. In an experimental discovery pipeline, candidates may not be independent draws from a well-defined repertoire distribution; selection from a trained model can introduce adaptive dependence. The sampling statement must then be re-derived or a deterministic hitting-set formulation used.

The theorem certifies intersection with thresholded single-agent ranges at a fixed assay context. It does not claim that selected antibodies work at a therapeutic dose in a combined product, that their interactions are non-antagonistic, or that each component is manufacturable and safe. Those require additional pharmacological and translational constraints.
\begin{figure}[t]\centering\includegraphics[width=.78\linewidth]{figures/fig04_sampling_success.pdf}\caption{Observed all-heavy-range hit fraction from 250 without-replacement samples per size in the synthetic finite model. These are finite-instance Monte Carlo diagnostics, not a real-data confidence guarantee.}\label{fig:mc}\end{figure}

\section{Interpreting the diversity-free claim}
Consider nested target panels $\Xcal_1\subseteq\Xcal_2\subseteq\cdots$. If a single range family $\Rcal_\infty$ contains every target-specific range, has VC dimension $d<\infty$, and every target range has $\mu$-mass at least $\eps$, the same theorem applies to every finite panel and to the entire declared family. This is stronger than selecting against one panel and hoping untested variants behave similarly: it is uniform over a mathematically defined target domain.

The conditions are demanding. If the infimum range mass over the domain is zero, no positive $\eps$ exists. A single rare but fully resistant isolate can invalidate the global assumption. One can restrict the domain, report a conditional certificate for a named region, stratify by clade or toxin family, or introduce a weighted exception set; one cannot silently retain a universal label. The domain restriction must be prospectively defined and clinically defensible.

A second issue is the meaning of $\mu$. Uniform weighting over a discovery library differs from frequency weighting over circulating strains, clinical exposure weighting, and cost weighting over manufacturable antibodies. A target range may have substantial mass under a laboratory library distribution but negligible mass after imposing developability or availability constraints. Such constraints belong in the candidate universe before $\mu$ is defined. If a measure is estimated from data, confidence bounds for range mass should be propagated into the certificate.

The formal content of “diversity-free” is therefore: panel cardinality alone need not force the theorem's sufficient bound to grow, provided range complexity and the lower-mass margin are controlled. It is not a theorem that biological diversity is irrelevant. Diversity can act precisely through $d$, $\eps$, response misspecification, and the definition of the target domain.

\section{Distinct-antibody coverage: existing epsilon-k-nets and a weighted bound}
Single coverage may be insufficient when an operational objective requires multiple distinct neutralizers. The corresponding combinatorial problem is not new: Alon, Jartoux, Keller, Smorodinsky, and Yuditsky studied the $\eps$-$t$-net problem as a generalization of classical $\eps$-nets \citep{AlonEtAl2022}. Their finite-hypergraph result constructs a family of $t$-element subsets such that every sufficiently large hyperedge contains at least one witness subset.

For a finite hypergraph $H=(V,\mathcal E)$ with $n=|V|$, VC dimension $d$, and polynomially bounded dual shatter function $\pi_H^*(m)\leq C m^{d^*}$, the cited theorem applies in its stated finite-size regime. For fixed $\eps\in(0,1)$ and integer $k\geq1$, it requires a sufficiently large candidate universe, with a condition of the form $n\geq C_1((k-1)/\eps)^{d^*}$, and yields pairwise-disjoint members of an $\eps$-$k$-net family with size
\begin{equation}|\mathcal S|=O\!\left(\frac{d(1+\log k)}{\eps}\log\frac{1}{\eps}\right),\label{eq:tknet}\end{equation}
where constants depend on dual-shatter growth. Finite VC dimension yields polynomial dual-shatter growth, but the induced exponent and pool-size threshold may be very large. This is a finite-uniform theorem, not a theorem for arbitrary continuous or weighted biological repertoires.

Since the family members are pairwise disjoint $k$-subsets, their union $U=\bigcup_{s\in\mathcal S}s$ has $|U|=k|\mathcal S|$ and intersects each heavy range at least $k$ times. Hence the conservative global union bound is
\begin{equation}|U|=O\!\left(\frac{k d(1+\log k)}{\eps}\log\frac{1}{\eps}\right).\label{eq:globalk}\end{equation}
This is an existing combinatorial route to distinct coverage under the finite-uniform hypotheses. It is not an immunological result and does not ensure epitope or escape-pathway independence.

For nonuniform measures, define the largest atom $\eta=\sup_{a\in\Acal}\mu(\{a\})$. Let $s$ be a uniform upper bound on the size of an $\eps/2$-net for every probability measure supported on a residual subset of candidates. The following elementary construction is a convenient sufficient condition.
\begin{lemma}[Weighted small-atom sequential $k$-cover]\label{lem:weighted} Let $(\Acal,\Rcal)$ have VC dimension $d<\infty$, let $\mu$ be a probability measure with maximum atom at most $\eta$, and suppose $\mu(R)\geq\eps$ for all $R\in\Rcal$. If $\eta\leq\eps/(2ks)$, then there exists $S\subseteq\Acal$ with $|S|\leq ks$ and $|S\cap R|\geq k$ for every range.\end{lemma}
\begin{proof} Select candidates in $k$ rounds. Before round $j+1$, at most $js$ candidates have been removed, so the removed mass is at most $js\eta\leq\eps/2$. Each original range therefore retains unnormalized mass at least $\eps/2$. After restriction and normalization, each range is still at least $\eps/2$-heavy. Choose an $\eps/2$-net of at most $s$ candidates on that residual measure. Repeat with unselected candidates. The rounds yield disjoint nets, so every range receives $k$ distinct intersections and the union has size at most $ks$.\end{proof}

The classical theorem at threshold $\eps/2$ gives a residual-net size of the order
\begin{equation}s=O\!\left(\frac{d}{\eps}\log\frac{2}{\eps}\right).\label{eq:residualnets}\end{equation}
Consequently, the sequential construction has total size
\begin{equation}|S|\leq ks=O\!\left(\frac{kd}{\eps}\log\frac{1}{\eps}\right).\label{eq:weightedsize}\end{equation}
This weighted lemma is a simple sufficient translation, not a new general multi-cover theorem and not claimed optimal. Its atom condition may be restrictive; finite weighted instances may be better handled by direct integer optimization.
\begin{figure}[t]\centering\includegraphics[width=.79\linewidth]{figures/fig06_multicover_cost.pdf}\caption{Greedy finite multi-cover size as demand $k$ increases. These are feasible synthetic solutions, not optimal cocktail sizes or values of the asymptotic theorem.}\label{fig:multicover}\end{figure}

\section{The price of redundancy and its biological interpretation}
The weighted sequential construction pays a factor of $k$ relative to the half-mass single-net size. This is a sufficient guarantee under its atom condition, not a lower bound on the true optimum. The finite-uniform $\eps$-$k$-net route pays a union-bound factor of $k$ and a further logarithmic dependence on $k$ in the cited theorem. Both approaches may be loose compared with an optimized finite incidence matrix.

The biological interpretation of redundancy must be explicit. Counting two antibodies is not equivalent to requiring two distinct epitopes, two independent escape pathways, or two agents with different pharmacokinetics. A stronger range definition can encode these properties, but its ground set and complexity may change. For example, define the candidate universe to be pairs $(a_1,a_2)$ satisfying an experimentally justified epitope-separation relation. A pair-level range is then a subset of antibody pairs, not simply a $k=2$ range in the original antibody space. The number of potential pairs grows quadratically in library size, and a formula for pair acceptance can increase VC complexity.

For HIV, distinctness could be defined using non-overlapping epitopes, distinct escape pathways, or resistance profiles. Each is a separate observable and requires structural or mutational evidence. For antivenom, redundancy could mean independent toxin-binding mechanisms or coverage of separate toxin families. It is not enough that two antibody molecules both bind a toxin if their functional effects are redundant or both fail against an unmeasured isoform.

General set multi-cover algorithms are established in computational geometry. Chekuri, Clarkson, and Har-Peled give approximation results for bounded-VC systems, including logarithmic approximation in the optimum for general bounded-VC settings and stronger factors in selected geometric cases \citep{Chekuri2009}. Those results address computational selection in finite instances. The present weighted lemma addresses existence under a mass-removal condition. Neither should be advertised as a new general multi-cover theory.
\begin{table}[t]\centering\small\caption{Coverage statements and the limits of their guarantees.}\begin{tabularx}{\linewidth}{p{.24\linewidth}p{.35\linewidth}X}\toprule Statement&Guaranteed property&Not guaranteed\\\midrule $\eps$-net&At least one selected candidate in each $\eps$-heavy range.&Exact optimum, combined-dose efficacy, out-of-domain protection.\\ Finite-uniform $\eps$-$k$-net&Union provides $k$ distinct hits under the theorem's hypergraph and size conditions.&Weighted or continuous candidate measures without additional work.\\ Weighted atom lemma&At least $k$ distinct candidates per heavy range under the atom condition.&Epitope independence, independent escape pathways, clinical redundancy.\\ Robust-neighborhood net&Hits every range defined by the specified perturbation neighborhood if its mass is at least $\eps$.&All future mutations or unknown toxin families.\\ Greedy finite multi-cover&A reproducible feasible cover for a finite incidence matrix.&Global optimality unless separately proved.\\\bottomrule\end{tabularx}\label{tab:guarantees}\end{table}

\section{Robust antigenic neighborhoods and escape}
A finite panel may conceal nearby variants that differ by unmeasured mutations or toxin composition. Let $E_\rho(x)\subseteq\Xcal$ be a declared neighborhood of target $x$, with radius or perturbation budget $\rho$. Define the robust range
\begin{equation}R_x^{\rm rob}(\rho)=\{a\in\Acal:\nu(a,z;c)\geq\theta\text{ for every }z\in E_\rho(x)\}.\label{eq:robustrange}\end{equation}
An antibody in this range covers the whole specified neighborhood, not just the nominal target. If the response relation and neighborhood are definable, the robust relation can often be expressed with quantifiers in the same o-minimal structure. First-order definability is closed under quantification, although the quantitative formula complexity can increase.

The robust margin is
\begin{equation}\eps_\rho=\inf_{x\in\Xcal}\mu(R_x^{\rm rob}(\rho)).\label{eq:robustmargin}\end{equation}
If neighborhoods are nested as $\rho$ increases, robust ranges can only shrink, so $\eps_\rho$ is non-increasing. A nominal $\eps$-net certificate does not persist automatically under a larger escape neighborhood. The relevant sufficient bound uses the robust margin itself, and may become vacuous if some robust range is empty.

A local neighborhood is not the full evolutionary future. Evolution can move outside a fixed radius, alter glycosylation or expression, change epitope accessibility, or create a new toxin family. “Escape-proof” should never be used without a bounded threat model and a defined update process. We therefore use \emph{neighborhood-robust} or \emph{escape-constrained}, not absolute escape-proof protection.
\begin{figure}[t]\centering\includegraphics[width=.80\linewidth]{figures/fig05_escape_neighbourhood.pdf}\caption{Synthetic robust candidate mass under four cardinal perturbations of increasing radius. The finite directions do not represent a continuum worst-case optimization.}\label{fig:escape}\end{figure}

\section{AI-assisted assay design: a bounded role}
The mathematical theorem does not require AI. AI becomes useful earlier, where the neutralization matrix is incomplete and assay capacity limited. A bounded role is uncertainty-aware experimental design: estimate which antibody--target pairs are near the classification boundary, prioritize measurements that most reduce uncertainty about the worst-case range mass, and update the certificate as new data arrive.

Let $Z_{a,x}$ be an unobserved assay outcome and $\mathcal D$ the current dataset. A probabilistic response model gives a posterior probability $p_{a,x}=\Prb(\nu(a,x;c)\geq\theta\mid\mathcal D)$. A cost-sensitive acquisition score can be
\begin{equation}\alpha(a,x)=\frac{\lambda_1\operatorname{Var}[\nu(a,x;c)\mid\mathcal D]+\lambda_2w_xp_{a,x}(1-p_{a,x})}{\operatorname{cost}(a,x)},\label{eq:acquisition}\end{equation}
where $w_x$ is a predeclared emphasis weight for target strata with uncertain or low estimated coverage. The first term values response-model uncertainty; the second focuses on uncertain threshold membership. This is a proposal, not a validated best policy.

A conservative certificate should use lower confidence bounds rather than posterior means alone. Let $\underline p_{a,x}$ be a calibrated lower bound for threshold membership. A candidate is counted only when $\underline p_{a,x}\geq1-\alpha_0$, where $\alpha_0$ is a predeclared assay risk tolerance. The lower mass estimate is
\begin{equation}\underline\mu(R_x)=\sum_iw_i\one\{\underline p_{i,x}\geq1-\alpha_0\},\qquad\sum_iw_i=1.\label{eq:empiricalmargin}\end{equation}
The system should issue a coverage certificate only if the lower bound exceeds $\eps$ for every declared target. If uncertainty is too large, it should abstain and request more assays. A model that confidently predicts broad coverage from sparse, miscalibrated data can create a false certificate; calibration, out-of-distribution testing, and prospective evaluation are therefore part of validation.

Conventional set-cover optimization can select the final components once an adequately measured incidence matrix exists. AI is justified only if it reduces the cost of acquiring the matrix or improves uncertainty control compared with ordinary random or uncertainty-based sampling. The active loop is shown in Fig.~\ref{fig:active}.
\begin{figure}[t]\centering\includegraphics[width=.83\linewidth]{figures/fig10_active_assay_loop.pdf}\caption{Proposed AI-assisted assay-selection loop. The model prioritizes measurements and updates lower confidence bounds; it must abstain when the coverage margin cannot be certified.}\label{fig:active}\end{figure}

\section{Finite-instance computational formulation}
For a finite library $A=\{a_1,\ldots,a_m\}$ and target panel $X=\{x_1,\ldots,x_n\}$, define an incidence matrix $C\in\{0,1\}^{n\times m}$ with $C_{ji}=1$ if candidate $a_i$ meets the prespecified endpoint for target $x_j$. The finite one-cover problem is
\begin{equation}\min_{z\in\{0,1\}^m}\sum_{i=1}^mz_i\quad\text{subject to}\quad\sum_{i=1}^mC_{ji}z_i\geq1\quad(j=1,\ldots,n).\label{eq:setcover}\end{equation}
The finite $k$-cover replaces the right side by $k$. Candidate costs can be added by minimizing $\sum_iw_iz_i$, where $w_i$ may represent manufacturing burden, dose burden, or another predeclared operational penalty. If interaction effects matter, a single-agent incidence matrix must be replaced by a combination-response model.

The VC theorem applies to an abstract range space and does not require that the finite optimization be solved exactly. Integer programming, greedy algorithms, LP rounding, or geometry-specific methods can be used. The greedy implementation below chooses the unused candidate that covers the largest number of currently unsatisfied target demands. It is deterministic under its tie-breaking rule and is not asserted to have a specific approximation ratio in this model.

Three outputs are kept separate: empirical candidate mass of each synthetic range; the random-sample frequency with which all heavy ranges are hit; and the greedy finite cover size for the realized incidence matrix. Empirical mass is a property of the finite synthetic pool; Monte Carlo frequency assesses that pool under without-replacement sampling; greedy cover size is a feasible integer-optimization solution. None evaluates the universal constant in the asymptotic theorem.
\begin{figure}[t]\centering\includegraphics[width=.86\linewidth]{figures/fig02_incidence_heatmap.pdf}\caption{A subset of the finite synthetic incidence matrix. Rows are antigen targets and columns are candidate antibodies; a filled entry means the 50\% Hill response threshold is met.}\label{fig:incidence}\end{figure}

\section{Synthetic model and verification protocol}
The demonstration uses a two-dimensional candidate feature plane. Candidate coordinates $z_i$ are sampled uniformly from $[0,1]^2$. Target centers $m_j$ are sampled from $[0.2,0.8]^2$, and target-specific radii $r_j$ are sampled from $[0.315,0.375]$. The IC$_{50}$ surface is set to $\IC(a_i,x_j)=\exp\{8(\|z_i-m_j\|_2-r_j)\}$, with Hill slope $h=1.4$ and concentration $c=1$. At threshold $\theta=0.5$, the response condition is exactly the disk condition $\|z_i-m_j\|_2\leq r_j$.

The fixed random seed is 20261101. The simulation has 4,000 candidate antibodies and 72 target variants. Candidate measure is empirical uniform mass over the finite pool; nominal threshold is $\eps=0.20$. This design creates a transparent range family whose VC dimension is at most three. It is not fitted to CATNAP, not inferred from a real antibody representation, and not intended to mimic a pathogen. It tests the algebraic equivalence between Hill threshold and geometric range, the finite multi-cover implementation, and monotonic loss of robust margin under increasing perturbations.

The random-sampling experiment draws subsets without replacement at sample sizes 8, 12, 16, 24, 32, 48, 64, 96, 128, and 192. Each sample size is evaluated over 250 trials. The finite greedy cover is solved for one and two hits per heavy range, then repeated for demands $k=1,2,3,4$. The robust-neighborhood experiment tests four cardinal perturbations at radii from 0 to 0.12. A candidate counts only if it covers all four shifted targets. This is a finite diagnostic, not a continuum worst-case optimization.
\begin{figure}[t]\centering\includegraphics[width=.78\linewidth]{figures/fig11_range_mass_distribution.pdf}\caption{Distribution of empirical candidate-repertoire masses for the 72 synthetic target ranges. Every target in this constructed instance exceeds the declared $\eps=0.20$ threshold.}\label{fig:mass}\end{figure}

\section{Numerical results}
The current deterministic run reports that all 72 synthetic target ranges are $\eps$-heavy. Their mean empirical mass is 0.359194, median mass is approximately 0.3614, and minimum mass is 0.256. A greedy one-cover selects three unique candidates; the greedy two-cover selects six unique candidates and achieves multiplicity between two and five across the target panel. These values describe the constructed finite incidence matrix only. They are not estimates of the number of antibodies required for HIV, snake venom, influenza, or any other biological target.

Across 250 repeated random-sampling trials per sample size, the observed fraction of samples hitting every heavy target range is 0.560 at size 8, 0.872 at size 12, 0.968 at size 16, 0.996 at size 24, and 1.000 at sizes 32 through 192. These are Monte Carlo frequencies under the synthetic design. They illustrate a transition in a finite toy instance; they do not estimate a universal theorem constant or provide an empirical confidence statement for a real candidate library.

Robust-neighborhood mass decreases with perturbation radius. In the current run, mean robust mass falls from 0.3592 at $\rho=0$ to approximately 0.1773 at $\rho=0.12$. The fraction of targets retaining at least $\eps=0.20$ mass decreases as well. This illustrates the formal point: a nominal coverage certificate does not persist automatically under a larger escape neighborhood. The robust margin must be evaluated for the robust ranges themselves.
\begin{table}[t]\centering\small\caption{Outputs from the deterministic synthetic verification run (seed 20261101).}\begin{tabularx}{\linewidth}{p{.49\linewidth}p{.20\linewidth}X}\toprule Metric&Observed value&Interpretation\\\midrule Candidate antibodies / targets&4,000 / 72&Synthetic feature-space instance\\ Mean / median / minimum range mass&0.3592 / 0.3614 / 0.2560&All 72 exceed $\eps=0.20$\\ Greedy one-cover / two-cover size&3 / 6&Feasible finite solutions; not proven optimal\\ Minimum multiplicity of two-cover&2&Every heavy target receives two distinct candidates\\ Hit fraction at $n=8$&0.560&250 without-replacement trials\\ Hit fraction at $n=24$&0.996&250 without-replacement trials\\ Hit fraction at $n=32$&1.000&250 without-replacement trials\\ Mean robust mass at $\rho=.12$&0.1773&Below nominal $\eps$ on average\\\bottomrule\end{tabularx}\label{tab:results}\end{table}
\begin{figure}[t]\centering\includegraphics[width=.79\linewidth]{figures/fig12_threshold_sensitivity.pdf}\caption{Sensitivity of range mass to stricter neutralization thresholds. A target can leave the heavy-range family as the response threshold increases.}\label{fig:threshold}\end{figure}

\section{Sensitivity, uncertainty, and model dependence}
The most consequential parameter is the lower-mass threshold $\eps$. The VC upper bound scales approximately as $\eps^{-1}\log(1/\eps)$, so a smaller worst-case neutralizer mass can increase sufficient selection size sharply. Threshold changes also alter range membership. Raising $\theta$ shrinks ranges and can reduce their masses. An empirical application should report a margin profile $\eps(\theta,c)$ across clinically relevant concentrations and functional endpoints rather than one convenient cut-off.

The candidate measure is equally important. Uniform weighting over a library, weighting by manufacturability, and weighting by exposure prevalence answer different questions. A rare high-potency class can be diluted by a large number of weak library members under a uniform measure. An availability-weighted measure may favor deployable agents but not represent immune repertoire probabilities. A measure based on exposure frequency can underweight rare but severe toxin profiles. Sensitivity analysis should vary plausible measures and show how the worst-case range mass changes.

Assay uncertainty is not binary. Replicates vary, neutralization curves can be censored near detection limits, and measurements across studies require harmonization. Treating each assay result as a deterministic edge can exaggerate precision. A stochastic alternative defines range membership through lower confidence bounds or explicitly models uncertain set membership. In either case, the coverage guarantee must account for uncertainty in both membership and the measure. A plug-in threshold set is exploratory unless uncertainty is propagated.

The two-dimensional disk model is especially favorable. High-dimensional sequence or structural feature spaces may be non-Euclidean and range boundaries may be nonconvex. Low VC dimension can still hold, but requires a separate derivation for the actual formula family. A compact learned embedding can produce smooth-looking regions while hiding sequence motifs that cause abrupt resistance. The range model should be validated on held-out target variants and study batches.
\begin{figure}[t]\centering\includegraphics[width=.78\linewidth]{figures/fig07_bipartite_graph.pdf}\caption{Sampled bipartite graph from the synthetic incidence matrix. Candidate antibody nodes and target nodes are joined by thresholded neutralization. This is a structural visualization, not a biological interaction network.}\label{fig:bipartite}\end{figure}

\section{Antibody cocktails are not vaccine immunogens}
The phrase “universal vaccine” can collapse two different design problems. A passive antibody cocktail administers selected neutralizing agents. Its ground set can be a library of antibodies, and the range definition closely matches the selection question, subject to pharmacokinetics and combination response constraints.

An active vaccine administers immunogens and relies on the host immune system to generate a distribution of antibody responses. Selecting immunogens is not selecting antibodies. A more faithful model introduces an induction operator
\begin{equation}\Psi:(I,h,s)\longmapsto\mu_{I,h,s},\label{eq:immunogenmap}\end{equation}
where $I$ is immunogen design, $h$ is host context, and $s$ is schedule/adjuvant context; $\mu_{I,h,s}$ is the induced antibody distribution or distribution over immune states. The margin then depends on immunogen design and host context. The $\eps$-net theorem can be applied conditionally to a realized induced repertoire only after $\mu_{I,h,s}$ is characterized. It does not give a direct finite bound on the number of immunogens required to induce that repertoire.

A vaccine-specific theory would couple antigen presentation, germinal-center dynamics, B-cell lineage selection, somatic hypermutation, memory formation, and serum concentration profiles to the induced measure. The operator $\Psi$ may be stochastic and not identifiable from a small set of serum measurements. Host heterogeneity, immune history, age, T-cell help, and immunodominance can change the induced distribution. These mechanisms are not represented by the present model.

Antivenom selection is closer to direct antibody selection, but binding does not guarantee functional neutralization. Toxin families may act through distinct pathways, venom composition varies by species and geography, and pharmacokinetic access matters. A range should be defined using a functional neutralization or rescue endpoint, not binding alone. The formalism is reusable only after target domain, endpoint, and measure are specified.
\begin{figure}[t]\centering\includegraphics[width=.78\linewidth]{figures/fig08_overlap_graph.pdf}\caption{Candidate-overlap graph for a synthetic subset. Edges indicate Jaccard overlap above a fixed threshold between target sets neutralized by two candidates. Overlap is not a proxy for epitope identity without structural evidence.}\label{fig:overlap}\end{figure}

\section{A translational protocol for HIV neutralization panels}
A first empirical study could use CATNAP as a source of published HIV neutralization measurements. The resource documents IC$_{50}$ and IC$_{80}$ values paired with virus sequence metadata \citep{Yoon2015}. The protocol should begin by fixing the endpoint, concentration convention, and treatment of censored values. The target set should be defined prospectively, with stratification by subtype, geography, study, and sequence quality where justified.

The candidate universe should contain antibodies with sufficiently harmonized measurements and stable identifiers. Missing assay pairs must remain missing, not be interpreted as failures or successes. Between-study variation and repeated measurements should be modeled, and the measure $\mu$ declared. A useful baseline is uniform mass over the antibody library; sensitivity analyses can consider availability weights or predeclared development constraints. Target ranges are the sets of antibodies meeting a fixed functional threshold for each virus, and empirical range mass is the weighted fraction of candidates in each range.

Three tasks are distinct. First, estimate the smallest observed range mass and uncertainty. Second, fit and test a restricted response model on held-out neutralization measurements. Third, estimate or bound the complexity of its range family. A feature map must be trained in separate folds or by cross-fitting to avoid using the same outcomes to make a range family artificially simple. The model class and complexity bound should be frozen before examining the final test panel.

If the lower confidence bound on the minimum range mass is zero or below the prespecified $\eps$, no positive-margin certificate is available for the full domain. The analysis may report a restricted conditional certificate for a subset of target strata, but every exclusion must be named. The selected net should be compared with direct greedy or integer-programming selection on the observed incidence matrix. This is a fair comparison: the theorem supplies an asymptotic sufficient bound, while optimization uses the observed finite panel and may find a much smaller set.

A toxin-focused study needs its own taxonomy, endpoint, and measure. Proteomic abundance, species prevalence, geographic distribution, functional toxin activity, and patient exposure are distinct weighting questions. A pooled toxin list can conceal a low-abundance but clinically decisive toxin class. Stratified margins or risk-weighted objectives may be more appropriate than a single global $\eps$.

\section{Falsifiable predictions and failure criteria}
The framework yields testable but conditional predictions. First, if an application has a fixed low-VC range family and a positive lower-mass margin, increasing the number of observed target variants without changing those quantities should not force the theoretical sufficient expression to grow linearly with panel size. This can be tested with nested panels and held-out target variants. It fails if inferred complexity grows, the minimum mass falls, or the model family changes.

Second, a fitted Hill/Pfaffian family should exhibit stable threshold geometry across held-out assay batches if it is an adequate response model. Systematic failures on unseen neutralization memberships, especially near glycan or epitope changes, undermine the application of that family. O-minimality of the mathematical model does not prove biological correctness.

Third, minimum range mass should be a stronger predictor of the existence of a small hitting set than raw target count once range complexity is controlled. This can be examined with matched synthetic and empirical panels. If cocktail size consistently scales with panel cardinality after controlling for margin and complexity, the model or complexity measure is missing structure.

Fourth, robust-neighborhood coverage should deteriorate as the threat neighborhood expands unless response breadth is stable under those perturbations. This can be falsified if robust range mass remains constant over a prospectively declared neighborhood. Fifth, active assay selection should reduce measurements required to certify a fixed margin compared with random or conventional uncertainty sampling, while preserving calibration and false-certificate control. A model that reduces assays by underestimating uncertainty has not succeeded.

\section{Competing explanations and confounders}
Apparent broad coverage can arise for reasons other than genuinely broad recognition. Candidate panels can overrepresent familiar strains, omit resistant lineages, or contain correlated laboratory isolates. Assay floor and ceiling effects can turn weak neutralization into apparent threshold membership. Study-to-study shifts in virus production, cell lines, reagent lots, or analysis conventions can change IC$_{50}$ estimates. A low-dimensional embedding may appear smooth while hiding sequence motifs that cause abrupt resistance.

Biological interaction effects are particularly important. Two antibodies can compete for overlapping epitopes, sterically interfere, or change effective concentrations. Conversely, a cocktail may exhibit synergy not present in single-agent assays. Fc-mediated function, tissue distribution, half-life, and timing of exposure influence protection independently of the single-agent neutralization matrix. The range-space certificate is therefore a screening statement, not a pharmacological prediction of a combined product.

For antivenoms, toxin abundance can create mismatch between the test panel and clinical mixture. A binder can recognize toxin without blocking toxic activity. Toxins may be present in low abundance but have high lethality, making uniform target weights inappropriate. A risk-weighted target measure should be justified by toxin activity and exposure evidence, not chosen to improve the margin.

A strong null explanation for a small finite cocktail is a favorable observed panel rather than a low-VC theorem. Exact integer optimization can find a small hitting set because the incidence matrix happens to have overlapping columns. A theoretical claim requires evidence that this behavior persists under held-out targets and a fixed model class. Report finite optimization and conditional VC bounds separately.

\section{Uncertainty-aware coverage certificates}
An empirical certificate should quantify uncertainty in both neutralization membership and the measure $\mu$. Suppose a finite library has replicate data and a model returns a simultaneous lower confidence bound $\underline p_{i,j}$ for candidate $i$ meeting target $j$'s threshold. A conservative operational range weight includes a candidate only when $\underline p_{i,j}\geq1-\alpha_0$, with predeclared assay risk tolerance $\alpha_0$. The lower mass estimate is given in Eq.~\eqref{eq:empiricalmargin}. Dependence across candidates and targets must be handled when simultaneous confidence is claimed.

A Bayesian construction can retain posterior membership probabilities and report a posterior distribution for each $\mu(R_j)$. One can require a lower posterior credible bound to exceed $\eps$ for all targets with a multiplicity adjustment or hierarchical model. This is model-dependent. Calibration on held-out assays and tests for distribution shift are needed. In small data regimes a conservative certificate may be impossible; that is a useful scientific outcome, not a reason to relax criteria after seeing the result.

The VC dimension is often unknown empirically. A practical strategy is to pre-specify a model family with a certified complexity bound and validate its predictive performance separately. Empirical shattering on a finite panel can provide a lower bound on population VC dimension, not an upper bound. Failure to find a shattered set of size $d+1$ in a finite dataset does not prove population VC dimension at most $d$.

\section{Research architecture and epistemic classification}
The claims fall into five categories. Hill-threshold equivalence, the classical $\eps$-net theorem, and the weighted sequential $k$-cover proof are P1 mathematical consequences under explicit assumptions. The use of a cost-sensitive acquisition score is a P2 design proposal. The hope that model complexity and margin remain stable as target panels grow is a P3-type statistical claim that is not established here. The mapping from antibody/antigen biology to a uniformly definable response family is a P4 structural correspondence pending empirical validation. The phrase “diversity-free” is P5 shorthand and is strictly qualified by fixed $d$ and positive $\eps$.

\begin{table}[t]\centering\small\caption{P1--P5 classification of central claims.}\begin{tabularx}{\linewidth}{p{.12\linewidth}p{.43\linewidth}X}\toprule Level&Claim&Status and evidence required\\\midrule P1&VC $\eps$-net bound; Hill threshold algebra; published $\eps$-$k$-net theorem; weighted small-atom lemma.&Exact within assumptions; proofs or prior theorem cited.\\ P2&Cost-aware assay acquisition and constrained cocktail design.&Explicit objectives; optimality depends on model and constraints.\\ P3&Range complexity remains stable as the target panel grows.&Not established; requires held-out, finite-size, and uncertainty analysis.\\ P4&Pfaffian/o-minimal response families model real antibody--antigen neutralization.&Structural correspondence proposed; empirical adequacy unresolved.\\ P5&“Diversity-free” shorthand for no explicit panel-size term.&Rhetorical summary only; qualified by fixed $d$ and positive margin.\\\bottomrule\end{tabularx}\label{tab:epistemic}\end{table}

A targeted search located classical $\eps$-nets, published $\eps$-$t$-nets, o-minimal and Pfaffian VC-dimension work, and general set-multicover algorithms. MathSciNet query execution did not return a verifiable result set in the available interface; zbMATH Open access was blocked. PubMed records were verified for the HIV bnAb reviews, recombinant antivenom review, and the 1 November 2026 human bnAb-lineage induction review. The novelty statement is narrow and provisional: this manuscript develops a particular conditional antibody-range formulation, explicit definability interface, and weighted sufficient-coverage translation. It does not claim first priority over differently named, unindexed, or subscription-only research.


\section{Minimax interpretation and what the upper bound cannot imply}
The range-space formulation can be interpreted as a one-sided robust selection problem. The designer chooses a finite set of candidates before the target range is revealed; the adversary chooses any target from the declared domain. The coverage objective is binary at a fixed assay threshold, and the adversary is constrained only by the lower-mass condition. The epsilon-net theorem therefore supplies a strategy for a particular information structure: the candidate set is sampled or selected from a known measure, while the target can be any member of the range family.

This is not equivalent to a clinical worst-case guarantee. A clinical adversary can change the target distribution, alter assay conditions, select a variant outside the declared domain, or exploit interactions among antibodies. The mathematical adversary in the theorem cannot do these things unless they are encoded by enlarging the family of ranges. If the target domain expands, the relevant range family and its VC dimension must be recomputed. If exposure probabilities shift, the measure must be updated and the minimum range mass re-estimated. The result is robust only against the choices allowed by the formal game.

There is a useful distinction between a distribution-free statement and a measure-free statement. VC epsilon nets are distribution-free in the sense that the theorem holds for an arbitrary fixed probability measure once every relevant range has measure at least epsilon. They are not measure-free: the value of epsilon and the set of heavy ranges depend on the chosen measure. Reweighting a repertoire after seeing the target panel can create an artificial margin and invalidate a prospective guarantee. Similarly, a finite library with uniform weights and a continuous repertoire with an atomless measure are different sampling spaces; converting one statement into the other requires a careful approximation argument.

Lower bounds are equally important. The upper bound of order $d\eps^{-1}\log(1/\eps)$ is not an exact prediction of the minimum number of antibodies. For some geometric range families, epsilon-net lower bounds show that logarithmic terms can be necessary; in others, special structure permits smaller sets. A biological response family may have additional constraints, such as monotone potency ordering or shared epitope classes, that enable a smaller optimum than the general worst-case bound. Conversely, correlated resistance or narrow epitope classes can produce difficult instances. The gap between the general upper bound and the observed minimum set is itself an empirical quantity worth studying.

A fair empirical evaluation should therefore report three quantities: (i) the general VC upper bound using a justified complexity certificate, (ii) the exact or approximate minimum of the finite incidence optimization, and (iii) a lower bound or hardness diagnostic for that finite instance. Reporting only the upper bound can make the framework appear less useful than it is; reporting only a tiny optimized finite set can obscure poor out-of-panel generalization. The minimax interpretation provides a common language for comparing these distinct quantities without conflating them.

\section{Identifiability of the measure, target domain, and response class}
The proposed certificate depends on three objects that are often estimated jointly: the target domain, the candidate measure, and the neutralization response class. This creates an identifiability problem. If the candidate library is assembled by an adaptive discovery process, its observed members are not necessarily representative of an underlying repertoire distribution. If the target panel is chosen because it is already known to be sensitive to available antibodies, the minimum range mass will be biased upward. If model architecture is selected after testing on the same panel, the apparent range complexity can be artificially reduced.

A pre-registered design can reduce, but not eliminate, these problems. Candidate identifiers and eligibility rules should be frozen before outcome analysis. The target domain should be declared independently of observed coverage, or a clear sampling mechanism should be recorded. Assay thresholds and missingness rules should be fixed prospectively. If the measure is estimated from a weighted library, uncertainty in the weights must be propagated. If the response function is learned, training and test targets must be separated at an appropriate biological level: random splitting of nearly identical isolates may overstate generalization.

There is also a distinction between the latent response surface and its thresholded range family. Different continuous response functions can induce the same range at a single threshold; conversely, small changes near the threshold can alter many memberships. VC dimension describes the collection of binary subsets, not the accuracy of the underlying concentration-response model. A low-VC threshold family can be systematically wrong. Therefore the empirical protocol must validate calibration and prediction at multiple clinically relevant thresholds in addition to complexity. The threshold itself is part of the scientific endpoint, not a harmless numerical parameter.

The minimum range mass is particularly sensitive to rare target strata. A pooled measure can hide a toxin family or viral clade that receives very little candidate mass. One option is stratified certification: specify target strata $X_1,\ldots,X_J$ and require a separate $\eps_j$ margin and a family-specific complexity $d_j$ in each. The resulting selected set may be a union of stratum-specific nets, with a size bounded by the sum of their guarantees. This can be more honest than one global certificate when the biology is heterogeneous, although it may cost additional candidates. A weighted risk objective is another option, but it changes the interpretation from universal coverage to risk-weighted coverage.

Finally, uncertainty in the VC dimension itself must not be hidden. Empirical shattering searches on finite observations can demonstrate that the true dimension is at least a certain value, but absence of an observed shattered configuration does not prove a population upper bound. A formal upper bound requires a model-family argument or a certified format count. If the application cannot supply such a certificate, the result should be reported as a finite-panel optimization study, not as a theoretical diversity-independent guarantee.

\section{Limitations}
The range-space model simplifies a high-dimensional biological process into thresholded neutralization sets. It abstracts immune development, antibody maturation, Fc effector mechanisms, pharmacokinetics, tissue access, toxicity, immunogenicity, manufacturing yield, epitope competition, cocktail synergy, viral dynamics, reservoirs, and evolutionary selection. The assumptions are not jointly validated by the synthetic model. Existence of a finite hitting set is not proof of in-vivo protection.

The qualitative o-minimal argument is robust for fixed definable families, but obtaining an explicit small VC dimension for a real neutralization model is a substantial future task. A complex sequence-to-response model may be definable but have a large bound. The paper does not derive a numerical Pfaffian format bound for HIV or toxins and does not set $d=2$ from antigenic cartography.

The multi-cover theorem requires a finite uniform hypergraph, dual-shatter growth, and sufficiently large candidate universe. The weighted small-atom lemma can require a large pool and can be loose. Neither combinatorial multiplicity statement guarantees independent epitopes or escape routes. The robust-neighborhood construction covers only the specified perturbation set.

The numerical study is entirely synthetic. The random-sampling results depend on one seed and one constructed finite geometry. The four-direction perturbation test is not a continuum worst-case optimization. The greedy cover is not certified optimal. No CATNAP values were downloaded and no antivenom assay data were analyzed. These limitations are explicit in the results files.

\section{Ethical, safety, and translational boundaries}
This manuscript is a mathematical research proposal, not a vaccine design recipe, antibody selection recommendation, antivenom formulation, or clinical protocol. The model does not identify a safe dose, establish efficacy, or support administration to humans or animals. Empirical follow-up must respect database reuse conditions and source-study provenance. Antibody discovery and animal testing require relevant institutional ethics and biosafety review; clinical translation requires toxicology, immunogenicity, pharmacokinetic, and regulatory evaluation.

A false claim of universal coverage could be harmful if used to justify reduced surveillance or underpowered clinical testing. A certificate must disclose target domain, assay context, candidate measure, threshold, lower confidence bound, model version, and excluded strata. A system unable to establish lower-mass margin should return “not certified” rather than a broad-spectrum label. No observed resistance in a finite panel is not a guarantee over an unobserved range family.

\section{Conclusions and next research steps}
We formulate broad-spectrum antibody selection as a range-space problem in which candidate antibodies are points and antigenic targets induce thresholded neutralization ranges. Under a positive lower-mass condition and finite VC dimension, classical $\eps$-net theory gives a sufficient selection bound of order $d\eps^{-1}\log(1/\eps)$ with no explicit dependence on target-panel cardinality. Uniform definability in an o-minimal structure provides a qualitative route to finite VC dimension; Pfaffian analysis can yield explicit bounds only after the response formula and format have been specified. Antigenic-map dimension is not the VC dimension of the antibody-range family.

For finite uniform candidate pools, the published $\eps$-$k$-net theorem gives a family of disjoint $k$-subsets whose union is a global $k$-fold cover under its finite hypergraph conditions. For weighted repertoires, we prove a separate sequential bound under a small-atom condition. Robust escape neighborhoods can be included by replacing nominal ranges with intersection ranges, but the robust margin may shrink and can vanish. The price of stronger robustness is therefore expressed through the margin and range complexity, not an unconditional cocktail-size constant.

The synthetic benchmark verifies a Hill-threshold disk family, finite random sampling, and greedy multi-cover for one disclosed toy instance. It does not validate biological response geometry. A next empirical step is reanalysis of a harmonized neutralization matrix such as CATNAP: fix candidate weights and endpoints, estimate uncertainty-aware range masses, fit and validate a constrained response model, and test whether a useful VC bound can be certified on held-out strains. A toxin-focused study needs a distinct target taxonomy and endpoints. A vaccine-specific theory must model the immunogen-to-induced-repertoire map rather than treating immunogens as antibodies.

The central conclusion is conditional: target count is not necessarily the only determinant of a sufficient antibody cocktail size. Lower mass of neutralization ranges, combinatorial complexity of the response family, and robustness margin under a declared escape neighborhood may be more informative. Whether these quantities can be estimated tightly enough to change real antibody design remains an empirical question.

\par\medskip
\noindent\begin{minipage}{\linewidth}\footnotesize\textbf{Research-process credit and disclosure.} Co-equal research-process credit is given to Devanik Debnath and Xylia, the autonomous multi-model research engine coordinating AI-assisted discovery, synthesis, and verification. Devanik Debnath retains responsibility for the final manuscript, scientific claims, citations, and released artifacts.\end{minipage}

\clearpage
\begin{thebibliography}{99}
\bibitem[Alon et~al.(2022)]{AlonEtAl2022} N. Alon, B. Jartoux, C. Keller, S. Smorodinsky, and Y. Yuditsky. The $\epsilon$-$t$-net problem. \emph{Discrete \& Computational Geometry}, 68(2):618--644, 2022. \href{https://doi.org/10.1007/s00454-022-00376-x}{doi:10.1007/s00454-022-00376-x}.
\bibitem[Chekuri et~al.(2009)]{Chekuri2009} C. Chekuri, K. L. Clarkson, and S. Har-Peled. On the set multi-cover problem in geometric settings. In \emph{Proceedings of the 25th Annual Symposium on Computational Geometry}, pp. 341--350, 2009. \href{https://doi.org/10.1145/1542362.1542421}{doi:10.1145/1542362.1542421}.
\bibitem[Haussler and Welzl(1987)]{HausslerWelzl1987} D. Haussler and E. Welzl. $\epsilon$-nets and simplex range queries. \emph{Discrete \& Computational Geometry}, 2(1):127--151, 1987. \href{https://doi.org/10.1007/BF02187876}{doi:10.1007/BF02187876}.
\bibitem[Haynes et~al.(2026)]{HaynesYehBaden2026} B. F. Haynes, C.-H. Yeh, and L. R. Baden. Induction of HIV broadly neutralizing antibodies with vaccination in humans in discovery medicine trials. \emph{Current Opinion in HIV and AIDS}, 21(6):442--449, 2026. \href{https://doi.org/10.1097/COH.0000000000001070}{doi:10.1097/COH.0000000000001070}.
\bibitem[Johnson and Laskowski(2010)]{JohnsonLaskowski2010} H. R. Johnson and M. C. Laskowski. Compression schemes, stable definable families, and o-minimal structures. \emph{Discrete \& Computational Geometry}, 43:914--926, 2010. \href{https://doi.org/10.1007/s00454-009-9201-3}{doi:10.1007/s00454-009-9201-3}.
\bibitem[Julg and Barouch(2021)]{JulgBarouch2021} B. Julg and D. Barouch. Broadly neutralizing antibodies for HIV-1 prevention and therapy. \emph{Seminars in Immunology}, 51:101475, 2021. \href{https://doi.org/10.1016/j.smim.2021.101475}{doi:10.1016/j.smim.2021.101475}.
\bibitem[Karpinski and Macintyre(1997)]{KarpinskiMacintyre1997} M. Karpinski and A. Macintyre. Polynomial bounds for VC dimension of sigmoidal and general Pfaffian neural networks. \emph{Journal of Computer and System Sciences}, 54(1):169--176, 1997. \href{https://doi.org/10.1006/jcss.1997.1477}{doi:10.1006/jcss.1997.1477}.
\bibitem[Menzies et~al.(2025)]{Menzies2025} S. K. Menzies, R. N. Patel, and S. Ainsworth. Practical progress towards the development of recombinant antivenoms for snakebite envenoming. \emph{Expert Opinion on Drug Discovery}, 20(6):799--819, 2025. \href{https://doi.org/10.1080/17460441.2025.2495943}{doi:10.1080/17460441.2025.2495943}.
\bibitem[Montana and Pardo(2009)]{MontanaPardo2009} J. L. Monta\~na and L. M. Pardo. On the Vapnik--Chervonenkis dimension of computer programs which use transcendental elementary operations. \emph{Annals of Mathematics and Artificial Intelligence}, 56:371--388, 2009. \href{https://doi.org/10.1007/s10472-009-9148-3}{doi:10.1007/s10472-009-9148-3}.
\bibitem[Sitaras(2020)]{Sitaras2020} I. Sitaras. Antigenic cartography: overview and current developments. In \emph{Influenza Virus: Methods and Protocols}, Methods in Molecular Biology, vol. 2123, pp. 61--68. Springer, 2020. \href{https://doi.org/10.1007/978-1-0716-0346-8_5}{doi:10.1007/978-1-0716-0346-8\_5}.
\bibitem[Thavarajah et~al.(2024)]{Thavarajah2024} J. J. Thavarajah, B. L. H\o nge, and C. M. Wejse. The use of broadly neutralizing antibodies (bNAbs) in HIV-1 treatment and prevention. \emph{Viruses}, 16(6):911, 2024. \href{https://doi.org/10.3390/v16060911}{doi:10.3390/v16060911}.
\bibitem[van den Dries et~al.(1994)]{vdDMM1994} L. van den Dries, A. Macintyre, and D. Marker. The elementary theory of restricted analytic fields with exponentiation. \emph{Annals of Mathematics}, 140(1):183--205, 1994.
\bibitem[Wilkie(1996)]{Wilkie1996} A. J. Wilkie. Model completeness results for expansions of the ordered field of real numbers by restricted Pfaffian functions and the exponential function. \emph{Journal of the American Mathematical Society}, 9(4):1051--1094, 1996.
\bibitem[Yoon et~al.(2015)]{Yoon2015} H. Yoon, J. Macke, A. P. West, B. Foley, P. J. Bjorkman, B. Korber, and K. Yusim. CATNAP: a tool to compile, analyze and tally neutralizing antibody panels. \emph{Nucleic Acids Research}, 43(W1):W213--W219, 2015. \href{https://doi.org/10.1093/nar/gkv404}{doi:10.1093/nar/gkv404}.
\end{thebibliography}

\appendix
\section{Derivation of the Hill threshold}
For positive Hill slope $h$, set $u=(c/\IC)^h$. Then $\nu=u/(1+u)$. The inequality $\nu\geq\theta$ is equivalent to $u\geq\theta/(1-\theta)$. Taking logarithms yields $h(\log c-\log\IC)\geq\log(\theta/(1-\theta))$, and division by $h$ gives Eq.~\eqref{eq:hillthreshold}. At $\theta=1/2$, the logarithmic ratio is zero and the condition reduces to $c\geq\IC$.

\section{Proof of the weighted small-atom lemma in expanded form}
Let $S_j$ be the union of the first $j$ residual nets, with $S_0=\varnothing$. Each round adds at most $s$ distinct points. For every range $R$, the mass removed before round $j+1$ is no larger than $|S_j|\eta\leq js\eta$. Under $\eta\leq\eps/(2ks)$ and $j\leq k-1$, the removed mass is at most $\eps/2$. Thus the unnormalized residual measure of $R$ is at least $\eps/2$. Normalizing the restricted measure divides by a total mass no greater than one, so its probability of $R$ is at least the unnormalized mass. A residual $\eps/2$-net therefore intersects $R$. Since each round samples only from the unselected candidate set, those intersections are distinct. After $k$ rounds the union covers each original range at least $k$ times and has size at most $ks$.

\section{Why the dual-shatter condition appears in the epsilon-k-net theorem}
The finite-uniform $\eps$-$k$-net theorem is formulated for a hypergraph and uses control of the dual shatter function. The dual family records which hyperedges contain a selected set of vertices. A polynomial bound on the number of distinct traces allows the theorem's counting and probabilistic arguments to control how many witness subsets are needed. Finite VC dimension implies polynomial growth of the dual shatter function by standard combinatorial bounds, but the corresponding dual exponent can be large. This is why finite VC dimension alone should not be paraphrased as a practical small-$n$ guarantee for every $k$.

\section{Reproducibility and data availability}
Run \verb|python verify_model.py|. The script requires Python 3.13, NumPy 2.3.5, Matplotlib 3.10.8, and NetworkX 3.6.1. It uses a fixed NumPy generator seed 20261101 and no network access, credentials, or external dataset. It generates the vector figures and synthetic CSV/JSON outputs in `figures/` and `results/`. No CATNAP table was downloaded and no antivenom assay dataset was analyzed in this release.
\end{document}